\originalchapter{第二次考前练习卷}
\psource{8.044，2013 春，Practice Exam 2，三题 35/30/35 分。中文题面；据公开官方答案核对的 AI 独立解答。}
\originalsection{1．弱相互作用玻色气体（35 分）}
\begin{exercise}
低温 $S=5aT^{3/2}V/2$、$\kappa_T=V^2/(2c)$，$a,c$ 恒定。低温、大体积极限中压力和能量密度趋零。
(a，15 分) 求 $P(T,V)$；(b，15 分) 求 $U(T,V)$；(c，5 分) 是否满足第三定律并解释。
\end{exercise}
\begin{solution}
(a) Maxwell 关系 $P_T=S_V=5aT^{3/2}/2$，故 $P=aT^{5/2}+g(V)$。压缩率给 $P_V=-2c/V^3$，所以 $P=aT^{5/2}+c/V^2+P_*$；边界要求 $P_*=0$。
(b) $U_T=TS_T=15aVT^{3/2}/4$，$U_V=TP_T-P=3aT^{5/2}/2-c/V^2$，积分得 $U=3aVT^{5/2}/2+c/V+U_*$。能量密度大体积极限不能固定有限加性常数 $U_*$，可选参考零点令其零。
(c) 固定 $V$，$S\to0$，符合此模型的第三定律行为。不存在把 $V\to\infty$ 与 $T\to0$ 任意联动而仍要求总熵趋零的普遍结论。
\end{solution}
\originalsection{2．一个有限热库（30 分）}
\begin{exercise}
可逆 Carnot 机在恒温无限冷库 $T_1$ 和有限恒容理想气体热库之间运行；热库初温 $T_2>T_1$、恒 $C_V$，直到降至 $T_1$。
(a，10 分) 热库熵变；(b，10 分) 机器对外功；(c，10 分) 全系统总熵变。
\end{exercise}
\begin{solution}
(a) $\Delta S_H=C_V\ln(T_1/T_2)$。(b) 可逆每步冷库得热 $T_1(-\dd S_H)$，总得热 $Q_C=T_1C_V\ln(T_2/T_1)$；热库放热 $Q_H=C_V(T_2-T_1)$，所以 $W_{\rm out}=C_V[T_2-T_1-T_1\ln(T_2/T_1)]>0$。(c) 冷库熵增抵热库熵降，循环机无熵净变，总零。
\end{solution}
\originalsection{3．三维经典超相对论气体（35 分）}
\begin{exercise}
体积 $V$ 中 $N$ 个不可分辨经典理想原子，温度足够高，$\epsilon=c|\vct p|$。
(a，7 分) 求正则 $Z$；(b，7 分) 求并画动量\textbf{大小}密度；(c，7 分) 内能；(d，7 分) 压力；(e，7 分) 熵。球坐标动量元 $p^2\sin\theta\dd p\dd\theta\dd\phi$。
\end{exercise}
\begin{solution}
(a) $z=8\pi V(\kb T)^3/(h^3c^3)$，$Z=z^N/N!$。(b) $p_r(p)=(\beta c)^3p^2e^{-\beta cp}/2$，峰在 $p=2\kb T/c$；不可漏掉球壳因子。(c) $U=3N\kb T$。(d) $P=\kb T\partial_V\ln Z=N\kb T/V$。(e) $S=\kb\ln Z+U/T=\kb[N\ln z-\ln N!+3N]\simeq N\kb[\ln(z/N)+4]$。适用还需稀薄、经典占据，不是温度高就可忽略一切相互作用。
\end{solution}
\begin{center}\begin{tikzpicture}\begin{axis}[width=7cm,height=3.6cm,xmin=0,xmax=8,ymin=0,ymax=.3,xlabel={$cp/(\kb T)$},ylabel={归一化密度},samples=100]\addplot[domain=0:8,thick]{x*x*exp(-x)/2};\end{axis}\end{tikzpicture}\end{center}
